\section{Final Results for PDF}\label{sec:final_results}

\begin{figure}[tbp]
\includegraphics[width=.48\textwidth]{images/unpolarized_different_momentum}
\caption{\label{fig:matchingPDFallPz} Nucleon boost momentum dependence of the matched unpolarized isovector PDFs: the dotted-green and solid-blue lines correspond to the nucleon momentum $P_z$ to be 1.8 and 2.3 GeV, respectively. The cutoff of Fourier transformation is chosen to be $zP^z\sim 12$ ($z=15a$ for $P_z=1.8$ GeV and $z=12a$ for $P_z=2.3$ GeV).}
\end{figure}

With the derivative method of FT and the matching using $p_z^R=0$ GeV and $\mu_R=3.2$ GeV, we show the dependence on the nucleon boosted momentum in Fig.~\ref{fig:matchingPDFallPz} with the statistical uncertainties.  They are consistent with each other as we can expect from the consistency of the quasi-PDF matrix element results in Fig.~\ref{fig:renormalized_h-tilde}.

 Finally, we show our results for PDF  and a comparison with global-analysis~\cite{Dulat:2015mca,Ball:2017nwa,Harland-Lang:2014zoa} in Fig.~\ref{fig:finalPDF}. As can be seen from the plot, our results show a reasonable agreement in the large-$x$ region, but at small-$x$ region there exists notable difference  majorly due to the systematic uncertainties from the FT truncation method and also the $p_z^R$ dependence.

\begin{figure}[tbp]
\includegraphics[width=.48\textwidth]{images/unpolarized_Pz=5_with_pheno}
\caption{\label{fig:finalPDF} Results for PDF at $\mu=2$~GeV calculated from RI/MOM quasi-PDF at nucleon momentum $P_z=2.3$~GeV: Comparing with CT14nnlo (90CL) \cite{Dulat:2015mca}, NNPDF3.1 (68CL) \cite{Ball:2017nwa}, and MMHT2014 (68CL) \cite{Harland-Lang:2014zoa}. Our results agree  with the global-analysis within   uncertainties.}
\end{figure}


